% Options for packages loaded elsewhere
\PassOptionsToPackage{unicode}{hyperref}
\PassOptionsToPackage{hyphens}{url}
\documentclass[
]{article}
\usepackage{xcolor}
\usepackage[margin=1in]{geometry}
\usepackage{amsmath,amssymb}
\setcounter{secnumdepth}{-\maxdimen} % remove section numbering
\usepackage{iftex}
\ifPDFTeX
  \usepackage[T1]{fontenc}
  \usepackage[utf8]{inputenc}
  \usepackage{textcomp} % provide euro and other symbols
\else % if luatex or xetex
  \usepackage{unicode-math} % this also loads fontspec
  \defaultfontfeatures{Scale=MatchLowercase}
  \defaultfontfeatures[\rmfamily]{Ligatures=TeX,Scale=1}
\fi
\usepackage{lmodern}
\ifPDFTeX\else
  % xetex/luatex font selection
\fi
% Use upquote if available, for straight quotes in verbatim environments
\IfFileExists{upquote.sty}{\usepackage{upquote}}{}
\IfFileExists{microtype.sty}{% use microtype if available
  \usepackage[]{microtype}
  \UseMicrotypeSet[protrusion]{basicmath} % disable protrusion for tt fonts
}{}
\makeatletter
\@ifundefined{KOMAClassName}{% if non-KOMA class
  \IfFileExists{parskip.sty}{%
    \usepackage{parskip}
  }{% else
    \setlength{\parindent}{0pt}
    \setlength{\parskip}{6pt plus 2pt minus 1pt}}
}{% if KOMA class
  \KOMAoptions{parskip=half}}
\makeatother
\usepackage{longtable,booktabs,array}
\usepackage{calc} % for calculating minipage widths
% Correct order of tables after \paragraph or \subparagraph
\usepackage{etoolbox}
\makeatletter
\patchcmd\longtable{\par}{\if@noskipsec\mbox{}\fi\par}{}{}
\makeatother
% Allow footnotes in longtable head/foot
\IfFileExists{footnotehyper.sty}{\usepackage{footnotehyper}}{\usepackage{footnote}}
\makesavenoteenv{longtable}
\usepackage{graphicx}
\makeatletter
\newsavebox\pandoc@box
\newcommand*\pandocbounded[1]{% scales image to fit in text height/width
  \sbox\pandoc@box{#1}%
  \Gscale@div\@tempa{\textheight}{\dimexpr\ht\pandoc@box+\dp\pandoc@box\relax}%
  \Gscale@div\@tempb{\linewidth}{\wd\pandoc@box}%
  \ifdim\@tempb\p@<\@tempa\p@\let\@tempa\@tempb\fi% select the smaller of both
  \ifdim\@tempa\p@<\p@\scalebox{\@tempa}{\usebox\pandoc@box}%
  \else\usebox{\pandoc@box}%
  \fi%
}
% Set default figure placement to htbp
\def\fps@figure{htbp}
\makeatother
\setlength{\emergencystretch}{3em} % prevent overfull lines
\providecommand{\tightlist}{%
  \setlength{\itemsep}{0pt}\setlength{\parskip}{0pt}}
\usepackage{bookmark}
\IfFileExists{xurl.sty}{\usepackage{xurl}}{} % add URL line breaks if available
\urlstyle{same}
\hypersetup{
  pdftitle={RN Baie de Saint-Brieuc},
  pdfauthor={Comptage ornithologique},
  hidelinks,
  pdfcreator={LaTeX via pandoc}}

\title{RN Baie de Saint-Brieuc}
\author{Comptage ornithologique}
\date{23 février 2026}

\begin{document}
\maketitle

\begin{figure}
\centering
\includegraphics[width=0.4\linewidth,height=\textheight,keepaspectratio]{logo/logo.png}
\caption{logo RN}
\end{figure}

\begin{center}\rule{0.5\linewidth}{0.5pt}\end{center}

\subsection{Comptage du 23/ 2 /2026}\label{comptage-du-23-2-2026}

\textbf{type de comptage :} Comptage Réserve Naturelle

\subsection{Conditions de comptage :}\label{conditions-de-comptage}

Ciel bleu/ Pas de précipitations /Mer calme

\begin{center}\rule{0.5\linewidth}{0.5pt}\end{center}

\subsection{Total d'oiseaux dénombrés lors du comptage
:}\label{total-doiseaux-duxe9nombruxe9s-lors-du-comptage}

\begin{longtable}[]{@{}lr@{}}
\toprule\noalign{}
& effectif \\
\midrule\noalign{}
\endhead
\bottomrule\noalign{}
\endlastfoot
Limicoles & 6133 \\
Anatidés & 1554 \\
autres & 634 \\
\end{longtable}

\begin{center}\rule{0.5\linewidth}{0.5pt}\end{center}

\newpage

\subsection{Détail du comptages :}\label{duxe9tail-du-comptages}

\begin{longtable}[]{@{}lrlll@{}}
\toprule\noalign{}
sp & effectif & precision & derangement & seuil \\
\midrule\noalign{}
\endhead
\bottomrule\noalign{}
\endlastfoot
Bécasseau variable & 1600 & - & - & - \\
Bécasseau sanderling & 1444 & - & - & National \\
Bécasseau maubèche & 1135 & - & - & National \\
Huîtrier pie & 901 & - & - & National \\
Courlis cendré & 883 & - & - & National \\
Bernache cravant & 808 & - & - & - \\
Grèbe huppé & 564 & - & - & National \\
Macreuse noire & 395 & - & - & National \\
Tadorne de Belon & 250 & - & - & - \\
Pluvier argenté & 139 & - & - & - \\
Canard colvert & 58 & - & - & - \\
Sarcelle d'hiver & 40 & - & - & - \\
Grand Cormoran & 28 & - & - & - \\
Héron garde-boeufs & 24 & - & - & - \\
Barge rousse & 15 & - & - & - \\
Tournepierre à collier & 11 & - & - & - \\
Aigrette garzette & 9 & - & - & - \\
Canard pilet & 3 & - & - & - \\
Chevalier aboyeur & 3 & - & - & - \\
Grèbe castagneux & 3 & - & - & - \\
Chevalier gambette & 2 & - & - & - \\
Oie cendrée & 2 & - & - & - \\
Pingouin torda & 2 & - & - & - \\
Eider à tête grise & 1 & - & - & - \\
Héron cendré & 1 & - & - & - \\
\end{longtable}

\newpage

\pandocbounded{\includegraphics[keepaspectratio]{C:/Users/Proprietaire/OneDrive - SAINT BRIEUC ARMOR AGGLOMERATION (1)/RNBaieSTBrieuc/SERENA/RserenaSQL/docs/ornitho-BilanComptagePDF_files/figure-latex/unnamed-chunk-3-1.pdf}}
\pandocbounded{\includegraphics[keepaspectratio]{C:/Users/Proprietaire/OneDrive - SAINT BRIEUC ARMOR AGGLOMERATION (1)/RNBaieSTBrieuc/SERENA/RserenaSQL/docs/ornitho-BilanComptagePDF_files/figure-latex/unnamed-chunk-3-2.pdf}}

\newpage

\subsection{Comparaison aux precedents
comptages}\label{comparaison-aux-precedents-comptages}

\begin{longtable}[]{@{}
  >{\raggedright\arraybackslash}p{(\linewidth - 8\tabcolsep) * \real{0.3433}}
  >{\raggedleft\arraybackslash}p{(\linewidth - 8\tabcolsep) * \real{0.1642}}
  >{\raggedright\arraybackslash}p{(\linewidth - 8\tabcolsep) * \real{0.1642}}
  >{\raggedright\arraybackslash}p{(\linewidth - 8\tabcolsep) * \real{0.1642}}
  >{\raggedright\arraybackslash}p{(\linewidth - 8\tabcolsep) * \real{0.1642}}@{}}
\toprule\noalign{}
\begin{minipage}[b]{\linewidth}\raggedright
sp
\end{minipage} & \begin{minipage}[b]{\linewidth}\raggedleft
23-02-2026
\end{minipage} & \begin{minipage}[b]{\linewidth}\raggedright
07-02-2026
\end{minipage} & \begin{minipage}[b]{\linewidth}\raggedright
24-01-2026
\end{minipage} & \begin{minipage}[b]{\linewidth}\raggedright
09-01-2026
\end{minipage} \\
\midrule\noalign{}
\endhead
\bottomrule\noalign{}
\endlastfoot
Aigrette garzette & 9 & 4 & - & 6 \\
Barge rousse & 15 & 23 & 66 & 2 \\
Bécasseau maubèche & 1135 & 1550 & 1655 & 1720 \\
Bécasseau sanderling & 1444 & 1794 & 1516 & 1996 \\
Bécasseau variable & 1600 & 2136 & 2275 & 2298 \\
Bernache cravant & 808 & 793 & 2263 & 966 \\
Canard colvert & 58 & 153 & 103 & 275 \\
Canard pilet & 3 & 97 & 100 & 78 \\
Chevalier aboyeur & 3 & 5 & 13 & 7 \\
Chevalier gambette & 2 & - & 90 & 2 \\
Courlis cendré & 883 & 873 & 476 & 711 \\
Eider à tête grise & 1 & 1 & - & 1 \\
Grand Cormoran & 28 & 22 & 41 & 16 \\
Grèbe castagneux & 3 & 3 & - & 5 \\
Grèbe huppé & 564 & 627 & 950 & 1155 \\
Héron cendré & 1 & 6 & - & 1 \\
Héron garde-boeufs & 24 & - & - & - \\
Huîtrier pie & 901 & 2022 & 2087 & 1238 \\
Macreuse noire & 395 & 753 & 380 & 428 \\
Oie cendrée & 2 & - & - & 2 \\
Pingouin torda & 2 & - & 32 & 7 \\
Pluvier argenté & 139 & 485 & 306 & 186 \\
Sarcelle d'hiver & 40 & 22 & 127 & 29 \\
Tadorne de Belon & 250 & 199 & 483 & 243 \\
Tournepierre à collier & 11 & 25 & 46 & 40 \\
\end{longtable}

\newpage

\subsection{Comparaison au mois de février
2025}\label{comparaison-au-mois-de-fuxe9vrier-2025}

\begin{longtable}[]{@{}lrll@{}}
\toprule\noalign{}
sp & 23-02-2026 & 04-02-2025 & 18-02-2025 \\
\midrule\noalign{}
\endhead
\bottomrule\noalign{}
\endlastfoot
Bécasseau variable & 1600 & 1751 & 1527 \\
Bécasseau sanderling & 1444 & 1882 & 1070 \\
Bécasseau maubèche & 1135 & 1085 & 900 \\
Huîtrier pie & 901 & 1514 & 2099 \\
Courlis cendré & 883 & 666 & 750 \\
Bernache cravant & 808 & 1006 & 724 \\
Grèbe huppé & 564 & 732 & 33 \\
Macreuse noire & 395 & 288 & 790 \\
Tadorne de Belon & 250 & 234 & 235 \\
Pluvier argenté & 139 & 224 & 225 \\
Canard colvert & 58 & 52 & 65 \\
Sarcelle d'hiver & 40 & 37 & 16 \\
Grand Cormoran & 28 & 31 & 29 \\
Héron garde-boeufs & 24 & - & - \\
Barge rousse & 15 & 79 & 82 \\
Tournepierre à collier & 11 & 25 & 13 \\
Aigrette garzette & 9 & 7 & 2 \\
Canard pilet & 3 & 47 & 93 \\
Chevalier aboyeur & 3 & 2 & 10 \\
Grèbe castagneux & 3 & 1 & 4 \\
Chevalier gambette & 2 & 116 & 45 \\
Oie cendrée & 2 & - & - \\
Pingouin torda & 2 & 80 & 37 \\
Eider à tête grise & 1 & 1 & 1 \\
Héron cendré & 1 & 1 & - \\
\end{longtable}

\newpage

\subsection{Synthèse des mois de février depuis
2000}\label{synthuxe8se-des-mois-de-fuxe9vrier-depuis-2000}

\textbf{\emph{Effectifs significativement supérieur à la moyenne}}

\begin{longtable}[]{@{}lrrlrr@{}}
\toprule\noalign{}
sp & 23-02-2026 & moy & ic & max & min \\
\midrule\noalign{}
\endhead
\bottomrule\noalign{}
\endlastfoot
Bécasseau sanderling & 1444 & 805.4 & (622-989) & 1882 & 9 \\
Courlis cendré & 883 & 644.5 & (565-724) & 1302 & 216 \\
Grèbe huppé & 564 & 267.0 & (150-384) & 1836 & 1 \\
Tadorne de Belon & 250 & 171.0 & (149-193) & 314 & 24 \\
Sarcelle d'hiver & 40 & 26.2 & (17-35) & 110 & 2 \\
Grand Cormoran & 28 & 19.4 & (16-23) & 46 & 3 \\
\end{longtable}

\emph{Effectifs non significativement supérieur à la moyenne}

\begin{longtable}[]{@{}lrrlrr@{}}
\toprule\noalign{}
sp & 23-02-2026 & moy & ic & max & min \\
\midrule\noalign{}
\endhead
\bottomrule\noalign{}
\endlastfoot
Héron garde-boeufs & 24 & 13.7 & (-8-35) & 24 & 1 \\
Aigrette garzette & 9 & 7.5 & (6-9) & 24 & 1 \\
Grèbe castagneux & 3 & 2.6 & (2-3) & 6 & 1 \\
Oie cendrée & 2 & 1.6 & (1-2) & 3 & 1 \\
Eider à tête grise & 1 & 1.0 & (1-1) & 1 & 1 \\
\end{longtable}

\emph{Effectifs non significativement inférieur à la moyenne}

\begin{longtable}[]{@{}lrrlrr@{}}
\toprule\noalign{}
sp & 23-02-2026 & moy & ic & max & min \\
\midrule\noalign{}
\endhead
\bottomrule\noalign{}
\endlastfoot
Macreuse noire & 395 & 452.1 & (320-584) & 1800 & 10 \\
Chevalier aboyeur & 3 & 5.2 & (3-7) & 15 & 1 \\
Héron cendré & 1 & 2.1 & (1-3) & 6 & 1 \\
\end{longtable}

\textbf{\emph{Effectifs significativement inférieur à la moyenne}}

\begin{longtable}[]{@{}lrrlrr@{}}
\toprule\noalign{}
sp & 23-02-2026 & moy & ic & max & min \\
\midrule\noalign{}
\endhead
\bottomrule\noalign{}
\endlastfoot
Bécasseau variable & 1600 & 1901.9 & (1663-2141) & 4250 & 180 \\
Bécasseau maubèche & 1135 & 1795.8 & (1522-2070) & 4750 & 330 \\
Huîtrier pie & 901 & 1785.6 & (1602-1970) & 3150 & 719 \\
Bernache cravant & 808 & 1145.7 & (893-1399) & 3420 & 253 \\
Pluvier argenté & 139 & 245.8 & (200-291) & 580 & 14 \\
Canard colvert & 58 & 111.7 & (87-136) & 347 & 29 \\
Barge rousse & 15 & 282.4 & (224-341) & 769 & 8 \\
Tournepierre à collier & 11 & 67.1 & (50-85) & 220 & 11 \\
Canard pilet & 3 & 75.9 & (58-93) & 230 & 2 \\
Chevalier gambette & 2 & 62.0 & (46-78) & 185 & 1 \\
Pingouin torda & 2 & 17.4 & (5-30) & 80 & 1 \\
\end{longtable}

\begin{center}\rule{0.5\linewidth}{0.5pt}\end{center}

\includegraphics[width=0.4\linewidth,height=\textheight,keepaspectratio]{logo/logo.png}
\pandocbounded{\includegraphics[keepaspectratio]{logo/Rserena.png}}
\pandocbounded{\includegraphics[keepaspectratio]{logo/RLogo.png}}
R-Serena /comptage/ version 3.0`

\end{document}
